\section{Bounce Perturbations and the Transfer Function}
\label{app:bounce}
%=============================================================================

\begin{tcolorbox}[colback=gray!5, colframe=gray!60, 
    title={\textbf{Status Note}}, fonttitle=\bfseries\small]
This appendix provides a minimal outline; the full numerical transfer-function evaluation
is implemented in the \texttt{bounce\_perturbations} verification module (see replication bundle).
\end{tcolorbox}

\subsection{L.1 Background Solution}
The bounce background is modeled by a smooth scale factor $a(\eta)$ where the torsion-induced
$a^{-6}$ term regularizes $z''/z$ for scalar perturbations. The transition to the $R^2$
inflationary phase is implemented by matching $a(\eta)$ and its derivatives at $\eta=\eta_t$.

\subsection{L.2 Numerical Integration Setup}
Mode functions are initialized in an adiabatic vacuum well before the bounce and integrated
through the bounce and into the $R^2$ era. The numerical scheme records the transfer function
by comparing late-time solutions to the pure $R^2$ baseline.

\subsection{L.3 Transfer Function}
Define the transfer function as
\begin{equation}
    T(k) \equiv \left|\frac{v_k(\eta\gg \eta_t)}{v_k^{(R^2)}(\eta\gg \eta_t)}\right|,
\end{equation}
with $T(k)\to 1$ for modes well inside the horizon during the bounce and deviations confined
to $k\sim k_{\text{bounce}}$.

\subsection{\texorpdfstring{L.4 The $k\to\ell$ Gate (CMB Observability as Policy Interface)}{L.4 The k->ell Gate (CMB Observability as Policy Interface)}}
\begin{tcolorbox}[colback=yellow!3, colframe=yellow!60!black,
    title={\textbf{What is deterministic vs.\ what is still policy}}, fonttitle=\bfseries\small]
\textbf{Deterministic (given the bounce completion):}
\begin{itemize}[leftmargin=1.5em, topsep=2pt, itemsep=1pt]
    \item The transfer-function shape $T(k)$ (this appendix; implemented in \texttt{bounce\_perturbations}).
    \item Given an explicit expansion-history budget, the mapping $k\to\ell$ and the resulting $C_\ell$ curves are deterministic (implemented in \texttt{k\_calibration} and \texttt{boltzmann\_transfer}).
\end{itemize}

\textbf{Policy / not yet ToE-closed:}
\begin{itemize}[leftmargin=1.5em, topsep=2pt, itemsep=1pt]
    \item The mapping from comoving $k$ to observed CMB multipole $\ell$ depends on the absolute expansion history (reheating/threshold bookkeeping), summarized as an overall scaling $a_0/a_t$ (or more precisely $a_0/a_{\rm transition}$ for the chosen transition surface).
    \item The signature classification (``CMB'' vs ``small-scale'') is therefore a \emph{gate}: it depends on the declared history, and should not be framed as an unconditional claim.
\end{itemize}

\textbf{Notation contract (to prevent $k$-confusion):}
\begin{itemize}[leftmargin=1.5em, topsep=2pt, itemsep=1pt]
    \item $k_{\hat{}}$ is a \emph{dimensionless} wave-number used inside the bounce solver (in $x=M\eta$ units).
    \item The raw bounce scales $k^{(s)}_{\rm bounce}$ and $k^{(t)}_{\rm bounce}$ are reported in these internal units; absolute $k(\mathrm{Mpc}^{-1})$ requires the TFPT scale $M$ and the expansion budget $a_0/a_{\rm transition}$.
    \item The observable proxy is $\ell \approx k(\mathrm{Mpc}^{-1})\,\chi_*$ with $\chi_*$ computed in a flat-$\Lambda$CDM snapshot (suite).
\end{itemize}

\textbf{Current suite snapshot (physics mode; threshold-derived budget):}
\begin{itemize}[leftmargin=1.5em, topsep=2pt, itemsep=1pt]
    \item $a_0/a_{\rm transition}\approx 2.4486\times 10^{57}$,
    \item $\chi_*\approx 1.3867\times 10^4$ Mpc,
    \item $\ell^{(s)}_{\rm bounce}\approx 56.997$ and $\ell^{(t)}_{\rm bounce}\approx 0.212$.
\end{itemize}
Under this explicit gate, the scalar bounce feature is CMB-low-$\ell$, while the tensor feature is super-horizon (below the quadrupole). 

\textbf{Falsifiable signature classes (conditional, gate-dependent):}
\begin{itemize}[leftmargin=1.5em, topsep=2pt, itemsep=1pt]
    \item \textbf{CMB-low-$\ell$ scalar feature}: features near $\ell\sim\mathcal{O}(10^2)$ (current default gate).
    \item \textbf{Small-scale-only}: features shifted to $\ell\gg 10^3$ under alternative histories.
\end{itemize}

\textbf{Suite evidence:} The \texttt{k\_calibration} module reports both (i) the required scaling to place features at target multipoles and (ii) the threshold-derived expansion-budget estimate; \texttt{boltzmann\_transfer} produces explicit CAMB-backed $C_\ell$ curves and the relative deviation $\Delta C_\ell/C_\ell$ between bounce-injected and baseline spectra.
\end{tcolorbox}

\TFPTkCalibrationTable

\begin{figure}[H]
\centering
\begin{tikzpicture}
\begin{groupplot}[
    group style={group size=2 by 1, horizontal sep=1.2cm},
    width=0.48\textwidth,
    height=0.42\textwidth,
    grid=both,
    legend style={font=\scriptsize},
]
\nextgroupplot[
    title={required $a_0/a_{\rm tr}$ vs target $\ell$},
    xmode=log,
    ymode=log,
    xlabel={target $\ell$},
    ylabel={$a_0/a_{\rm tr}$},
    xmin=1.5, xmax=1000,
]
\addplot+[mark=*] coordinates {
    (2,6.978134315108868e58)
    (30,4.652089543405912e57)
    (700,1.9937526614596766e56)
};
\addlegendentry{scalar (required)}
\addplot+[mark=square*] coordinates {
    (2,2.5901704412653403e56)
    (30,1.7267802941768936e55)
    (700,7.400486975043829e53)
};
\addlegendentry{tensor (required)}
\addplot[red, dashed, thick] coordinates {(2,2.4485807685210785e57) (700,2.4485807685210785e57)};
\addlegendentry{current budget ($2.45\times10^{57}$)}

\nextgroupplot[
    title={gate view (CMB window + bounce $\ell$)},
    xmode=log,
    xlabel={$\ell$},
    ymin=0, ymax=1,
    ytick=\empty,
    xmin=0.1, xmax=5000,
    axis y line=none,
]
% CMB window shading: [2,2500]
\addplot[draw=none, fill=blue!10] coordinates {(2,0) (2500,0) (2500,1) (2,1)} \closedcycle;
\addplot[black!40] coordinates {(2,0) (2,1)};
\addplot[black!40] coordinates {(2500,0) (2500,1)};
\node[anchor=south, font=\scriptsize] at (axis cs:20,0.02) {CMB window};

% bounce locations (physics-suite snapshot)
\addplot[green!60!black, thick] coordinates {(56.99737909257206,0) (56.99737909257206,1)};
\node[anchor=south, font=\scriptsize, green!60!black] at (axis cs:56.99737909257206,0.02) {$\ell_b^{(s)}\approx57$};
\addplot[orange!80!black, thick] coordinates {(0.2115650400072186,0) (0.2115650400072186,1)};
\node[anchor=south, font=\scriptsize, orange!80!black] at (axis cs:0.2115650400072186,0.02) {$\ell_b^{(t)}\approx0.21$};
\end{groupplot}
\end{tikzpicture}
\caption{Suite plots from \texttt{k\_calibration} (physics mode): scaling and feasibility diagnostics for the $k\to\ell$ gate under an explicit expansion-history budget.}
\end{figure}

\begin{figure}[H]
\centering
\begin{tikzpicture}
\begin{groupplot}[
    group style={group size=2 by 2, horizontal sep=1.2cm, vertical sep=1.0cm},
    width=0.48\textwidth,
    height=0.33\textwidth,
    xmode=log,
    xmin=2, xmax=3000,
    grid=major,
    minor tick num=0,
    ticklabel style={font=\scriptsize},
    legend style={font=\scriptsize},
]
\nextgroupplot[
    title={$D_\ell^{TT}$ (baseline vs bounce)},
    ymode=log,
    ylabel={$D_\ell^{TT}$ [$\mu$K$^2$]},
]
\addplot+[mark=*, mark size=2pt] coordinates {
    (2,1053.602432)
    (3,998.0887611)
    (4,944.9109792)
    (5,905.3057583)
    (10,845.7614033)
    (20,936.246946)
    (30,1090.649646)
    (50,1468.079395)
    (80,2185.897476)
    (100,2782.530282)
    (150,4517.548558)
    (200,5768.991349)
    (300,4189.951412)
    (500,2521.115802)
    (700,1928.629066)
    (1000,1093.569344)
    (1500,706.1670601)
    (2000,237.9250741)
    (2508,78.09810466)
};
\addlegendentry{baseline}
\addplot+[mark=square*, mark size=2pt] coordinates {
    (2,10994.05279)
    (3,10486.33276)
    (4,9385.765294)
    (5,7182.088179)
    (10,466.5297095)
    (20,311.2303573)
    (30,647.8609341)
    (50,1396.409529)
    (80,2178.001377)
    (100,2764.125993)
    (150,4513.504899)
    (200,5742.64759)
    (300,4165.032244)
    (500,2519.57907)
    (700,1945.981149)
    (1000,1115.817699)
    (1500,728.4072936)
    (2000,247.883187)
    (2508,81.98087433)
};
\addlegendentry{bounce injected}

\nextgroupplot[
    title={$D_\ell^{EE}$ (baseline vs bounce)},
    ymode=log,
    ylabel={$D_\ell^{EE}$ [$\mu$K$^2$]},
]
\addplot+[mark=*, mark size=2pt] coordinates {
    (2,0.03176782281)
    (3,0.0407463901)
    (4,0.03535540342)
    (5,0.02363620444)
    (10,0.003157610827)
    (20,0.005394191594)
    (30,0.0223578033)
    (50,0.1235977312)
    (80,0.4863197226)
    (100,0.8033626171)
    (150,1.135920327)
    (200,0.710593948)
    (300,9.377843402)
    (500,8.525653848)
    (700,38.4094504)
    (1000,43.34908241)
    (1500,13.22941615)
    (2000,8.84123834)
    (2508,3.009426377)
};
\addlegendentry{baseline}
\addplot+[mark=square*, mark size=2pt] coordinates {
    (2,0.5094182832)
    (3,0.562182037)
    (4,0.3663058727)
    (5,0.1665671877)
    (10,0.0005850428195)
    (20,0.0006453541167)
    (30,0.004809524067)
    (50,0.09804048815)
    (80,0.489652682)
    (100,0.7951178048)
    (150,1.13708347)
    (200,0.7069214332)
    (300,9.311700761)
    (500,8.492435159)
    (700,38.50920372)
    (1000,43.9978482)
    (1500,13.62199765)
    (2000,9.189186049)
    (2508,3.154751105)
};
\addlegendentry{bounce injected}

\nextgroupplot[
    title={relative deviation (TT)},
    title style={at={(0.5,0.97)}, yshift=-0.6ex},
    xlabel={$\ell$},
    ylabel={$\Delta D_\ell/D_{\ell,0}$},
    ymin=-1.0, ymax=10.5,
]
\addplot+[mark=*, mark size=2pt] coordinates {
    (2,9.43472609)
    (3,9.506413022)
    (4,8.93296247)
    (5,6.933328727)
    (10,-0.4483908727)
    (20,-0.6675766381)
    (30,-0.405986206)
    (50,-0.0488187938)
    (80,-0.003612291319)
    (100,-0.006614227665)
    (150,-0.0008951004397)
    (200,-0.004566441139)
    (300,-0.005947364325)
    (500,-0.0006095444869)
    (700,0.008997107253)
    (1000,0.02034471301)
    (1500,0.03149429467)
    (2000,0.04185398693)
    (2508,0.04971656719)
};
\addplot[black!50, dashed] coordinates {(2,0) (2508,0)};

\nextgroupplot[
    title={relative deviation (EE)},
    title style={at={(0.5,0.97)}, yshift=-0.6ex},
    xlabel={$\ell$},
    ylabel={$\Delta D_\ell/D_{\ell,0}$},
    ymin=-1.0, ymax=16.5,
]
\addplot+[mark=*, mark size=2pt] coordinates {
    (2,15.03566874)
    (3,12.79710045)
    (4,9.360675802)
    (5,6.04712079)
    (10,-0.8147197829)
    (20,-0.8803612913)
    (30,-0.7848838724)
    (50,-0.2067776069)
    (80,0.006853432569)
    (100,-0.01026287768)
    (150,0.00102396489)
    (200,-0.005168232629)
    (300,-0.007053075852)
    (500,-0.003896321651)
    (700,0.002597103541)
    (1000,0.0149660788)
    (1500,0.02967489204)
    (2000,0.0393550876)
    (2508,0.04828984317)
};
\addplot[black!50, dashed] coordinates {(2,0) (2508,0)};
\end{groupplot}
\end{tikzpicture}
\caption{Suite plots from \texttt{boltzmann\_transfer} (physics mode): CAMB-backed $C_\ell$ curves (TT/EE) for baseline vs bounce-injected primordial spectrum, and the relative deviation $\Delta C_\ell/C_\ell$.}
\end{figure}

\subsection{L.5 Falsification Targets}
Large-scale suppression or oscillatory features in $\mathcal{P}_\zeta(k)$ are predicted at
$k\lesssim k_{\text{bounce}}$. Current CMB large-scale bounds constrain the allowable feature
amplitude, providing a direct observational test of the bounce completion.

%=============================================================================
\section{BH1 Dyonic Toy Solution and Polarization Rotation Scaling (Appendix)}
\label{app:bh-dyonic}
%=============================================================================

\begin{tcolorbox}[colback=yellow!3, colframe=yellow!60!black, breakable,
title={\textbf{Scope (toy model; not a kernel claim)}}, fonttitle=\bfseries\small]
This appendix records a minimal dyonic toy solution showing the $1/r^2$ scaling of the axion profile and the induced polarization rotation.
It is included as analytical support for the EHT intercept test (Sec.~\ref{sec:eht-intercept}), not as a standalone astrophysical claim.
\end{tcolorbox}

\textbf{Axion equation sourced by $F\tilde F$.}
Varying the UFE action (Definition~\ref{def:UFEaction}) gives, in a flat static approximation and neglecting $V'(a)$,
\begin{equation}
\nabla^2 a \simeq \cthree\, F_{\mu\nu}\tilde F^{\mu\nu} = -4\cthree\,\mathbf{E}\cdot\mathbf{B}.
\end{equation}

\textbf{Dyonic fields.}
For a pointlike dyon with effective charges $(Q_e,Q_m)$,
\begin{equation}
\mathbf{E}(r)=\frac{Q_e}{4\pi}\frac{\hat r}{r^2},\qquad
\mathbf{B}(r)=\frac{Q_m}{4\pi}\frac{\hat r}{r^2}
\quad\Rightarrow\quad
\mathbf{E}\cdot\mathbf{B}=\frac{Q_eQ_m}{16\pi^2}\frac{1}{r^4}.
\end{equation}
Using $\nabla^2 (1/r^2)=2/r^4$ for $r\neq 0$, a particular solution is
\begin{equation}
a(r)=a_\infty+\frac{\cthree Q_e Q_m}{8\pi^2}\frac{1}{r^2}.
\end{equation}

\textbf{Polarization rotation scaling.}
The birefringence transport law (Theorem~\ref{thm:birefringence}) implies $\beta = 2\cthree\,\Delta a$ for an accumulated axion excursion $\Delta a$.
For emission at radius $r_{\mathrm{em}}$ observed at infinity,
\begin{equation}
\Delta a = a_\infty - a(r_{\mathrm{em}}) = -\frac{\cthree Q_e Q_m}{8\pi^2 r_{\mathrm{em}}^2}
\quad\Rightarrow\quad
\beta_{\mathrm{BH}}(r_{\mathrm{em}})=\frac{Q_e Q_m}{256\pi^4 r_{\mathrm{em}}^2},
\end{equation}
matching the scaling used in Sec.~\ref{sec:eht-intercept} (up to sign conventions).\hfill\cite{blackHolesBh1Note2025}

%=============================================================================
